%Abstract
\setlength{\headheight}{14pt}
\begin{center}
	{\huge \bf Abstract}
	\line(1,0){430}
\end{center}

Student dropout prediction is a critical challenge in higher education, with profound implications for institutional planning, resource allocation, and student success. Traditional machine learning approaches often overlook temporal patterns and psychosocial factors that influence dropout behavior. This thesis presents a comprehensive investigation of student dropout prediction using a multi-tiered methodology combining classical machine learning, ensemble methods, deep learning, and a novel hybrid framework integrating temporal attention mechanisms with Large Language Model (LLM)-derived psychosocial features. Analyzing 4,424 students across 46 features (18 academic, 12 financial, 16 demographic), we systematically evaluate six baseline models (Decision Tree, Naive Bayes, Random Forest, AdaBoost, XGBoost, Neural Network) using rigorous 10-fold cross-validation. Feature ranking analysis using Information Gain, Gain Ratio, and Gini Index identifies Curricular Units (2nd semester approved), Tuition Fees status, and semester grades as the most influential dropout predictors. The proposed Adaptive Hierarchical Feature Selection with Temporal Attention (AHFS-TA) framework achieves state-of-the-art performance with 91.32\% accuracy and 95.5\% AUC-ROC, outperforming all baselines. Explainable AI techniques (SHAP values, attention weights) provide interpretable insights for early intervention. The research demonstrates that multimodal learning combining structured educational data with LLM-extracted features (+1.71\% improvement) and temporal progression modeling (+1.18\% gain) significantly enhances prediction accuracy. Results validate adaptive feature selection reduces dimensionality by 26\% while improving performance. This work contributes: (1) comprehensive empirical evaluation of six models on a large educational dataset, (2) rigorous feature ranking analysis identifying critical dropout indicators, (3) novel AHFS-TA framework integrating LLM features with temporal attention, achieving 91.32\% accuracy and exceeding published benchmarks by +4.02\%, and (4) explainable AI framework providing actionable insights for educational stakeholders. Findings enable proactive intervention strategies, optimized resource allocation, and personalized student support systems.
\clearpage
\setlength{\headheight}{12pt}